% Fuente única de clase: Recursos/2025-10-16.pdf, páginas 1--10.
En la clase 17 calculamos máximos comunes divisores mediante el algoritmo de Euclides. Después usamos el lema de Zorn para obtener ideales maximales y estudiamos primitividad y contenido de polinomios. El lema de Gauss preparará la factorización en anillos de polinomios.

%%%%%%%%%%%%%%%%%%%%%%%%%%%%%%%%%%%%%%%%%%%%%%%%%%%%%%%%%%%%%%%%%%%%%%%%%%%%%%%
\section{El algoritmo de Euclides}
%%%%%%%%%%%%%%%%%%%%%%%%%%%%%%%%%%%%%%%%%%%%%%%%%%%%%%%%%%%%%%%%%%%%%%%%%%%%%%%
%%%%%%%%%%%%%%%%%%%%%%%%%%%%%%%%%%%%%%%%%%%%%%%%%%%%%%%%%%%%%%%%%%%%%%%%%%%%%%%
\subsection{División y conservación de los divisores comunes}
%%%%%%%%%%%%%%%%%%%%%%%%%%%%%%%%%%%%%%%%%%%%%%%%%%%%%%%%%%%%%%%%%%%%%%%%%%%%%%%
% P01--P02.
Recordemos que un dominio euclidiano admite una función $v:R\setminus\{0\}\to\mathbb Z_{\geq0}$ para la cual toda división por $b\ne0$ tiene la forma
\[
a=bq+r,\qquad r=0\quad\text{o}\quad v(r)<v(b).
\]
No se compara directamente $r$ con $b$ mediante un orden del anillo: se comparan sus tamaños. Como se probó en la clase 16, estos dominios son DIP: un elemento no nulo de tamaño mínimo en un ideal lo genera, porque todo resto en ese ideal debe ser cero. Esto incluye $\mathbb Z$, con $v(n)=|n|$, y $K[x]$, con $v(f)=\deg f$, para un campo $K$.

\begin{lema}\label{lema:C17-resto}
En un anillo conmutativo, si $a=bq+r$, entonces $(a,b)=(b,r)$.
\end{lema}
\begin{proof}
Como $r=a-bq$, los generadores $b,r$ pertenecen a $(a,b)$. Recíprocamente, $a=bq+r$ pertenece a $(b,r)$ y también $b$. Esto prueba las dos inclusiones.
\end{proof}
\begin{coro}\label{coro:C17-mcd-resto}
En un dominio, los pares $a,b$ y $b,r$ tienen los mismos divisores comunes. En particular, uno tiene máximo común divisor si y sólo si lo tiene el otro, y sus máximos comunes divisores son asociados.
\end{coro}
\begin{proof}
Un elemento $c$ divide a ambos generadores de un ideal si y sólo si ese ideal está contenido en $(c)$. Por el lema, esta condición es la misma para los dos pares. También lo es, por tanto, la propiedad que define al máximo común divisor.
\end{proof}
Las igualdades entre máximos comunes divisores se entienden salvo asociados, o tras elegir representantes compatibles.

%%%%%%%%%%%%%%%%%%%%%%%%%%%%%%%%%%%%%%%%%%%%%%%%%%%%%%%%%%%%%%%%%%%%%%%%%%%%%%%
\subsection{Terminación y cálculo del último resto}
%%%%%%%%%%%%%%%%%%%%%%%%%%%%%%%%%%%%%%%%%%%%%%%%%%%%%%%%%%%%%%%%%%%%%%%%%%%%%%%
% P02--P03.
Sean $a,b\in R$, con $b\ne0$, y pongamos $r_{-1}=a$, $r_0=b$. Mientras $r_{i-1}\ne0$, hacemos una división
\[
r_{i-2}=q_i r_{i-1}+r_i.
\]
\begin{prop}[Algoritmo de Euclides]\label{prop:C17-Euclides}
El procedimiento termina con un resto cero. Si $r_N=0$ es el primero, entonces $r_{N-1}$ es un máximo común divisor de $a,b$.
\end{prop}
\begin{proof}
Si nunca apareciera cero, tendríamos una sucesión estrictamente decreciente
\[
v(r_0)>v(r_1)>v(r_2)>\cdots
\]
de enteros no negativos. Es imposible, así que existe ese $N$. El último paso dice $r_{N-2}=q_Nr_{N-1}$, de donde $(r_{N-2},r_{N-1})=(r_{N-1})$. Aplicando el \cref{lema:C17-resto} a cada división obtenemos
\[
(a,b)=(r_0,r_1)=\cdots=(r_{N-2},r_{N-1})=(r_{N-1}).
\]
Su generador es un máximo común divisor. Si $N=1$, la misma conclusión proviene directamente de $a=q_1b$.
\end{proof}
Si $b=0$, un máximo común divisor es $a$; si ambos elementos son cero, se puede tomar cero. En un campo toda división por un elemento no nulo tiene resto cero, por lo que cualquier función de tamaño cumple la condición euclidiana.

\begin{ejem}
En $\mathbb Q[x]$, sean
\[
a=x^3-x^2-4x+4,\qquad b=x^3+x^2-4x-4.
\]
Las divisiones son
\[
\begin{aligned}
a&=b+(-2x^2+8),\\
b&=(-2x^2+8)\left(-\frac{x}{2}-\frac12\right).
\end{aligned}
\]
El último resto no nulo es $-2x^2+8$. Por tanto, es un máximo común divisor. Su asociado mónico es $x^2-4$, que también sirve como representante.
\end{ejem}

%%%%%%%%%%%%%%%%%%%%%%%%%%%%%%%%%%%%%%%%%%%%%%%%%%%%%%%%%%%%%%%%%%%%%%%%%%%%%%%
\section{Orden, elección e ideales maximales}
%%%%%%%%%%%%%%%%%%%%%%%%%%%%%%%%%%%%%%%%%%%%%%%%%%%%%%%%%%%%%%%%%%%%%%%%%%%%%%%
%%%%%%%%%%%%%%%%%%%%%%%%%%%%%%%%%%%%%%%%%%%%%%%%%%%%%%%%%%%%%%%%%%%%%%%%%%%%%%%
\subsection{El lema de Zorn y el buen orden}
%%%%%%%%%%%%%%%%%%%%%%%%%%%%%%%%%%%%%%%%%%%%%%%%%%%%%%%%%%%%%%%%%%%%%%%%%%%%%%%
% P04--P05: enunciados fundacionales; las pizarras no prueban su equivalencia.
\begin{defi}
Un \emph{orden parcial} en un conjunto $Z$ es una relación $\preceq$ reflexiva, transitiva y antisimétrica: $a\preceq b$ y $b\preceq a$ implican $a=b$. Es \emph{total} si cualesquiera dos elementos son comparables. Una \emph{cadena} es un subconjunto totalmente ordenado por la relación restringida.
\end{defi}
El orden usual de $\mathbb Z$ es total. La inclusión en una familia de subconjuntos es un orden parcial y, en general, no es total. Escribimos $a\prec b$ si $a\preceq b$ y $a\ne b$.

Un elemento $m$ es \emph{maximal} si $m\preceq a$ implica $m=a$. Una \emph{cota superior} de $S\subseteq Z$ es un elemento $u\in Z$ con $a\preceq u$ para todo $a\in S$. Ser maximal no exige estar por encima de todos los elementos: puede haber elementos incomparables.

\begin{lema}[Zorn]\label{lema:C17-Zorn}
Sea $Z$ un conjunto no vacío parcialmente ordenado. Si toda cadena tiene una cota superior en $Z$, entonces $Z$ tiene un elemento maximal.
\end{lema}
\begin{defi}
Un orden total es un \emph{buen orden} si todo subconjunto no vacío tiene un elemento mínimo. Un conjunto es \emph{bien ordenable} si admite un buen orden.
\end{defi}
\begin{teo}[Buen orden]
Todo conjunto admite un buen orden.
\end{teo}
El \emph{axioma de elección} permite seleccionar un elemento de cada miembro de una familia de conjuntos no vacíos; también puede formularse para familias de conjuntos disjuntos. Sobre los axiomas de Zermelo--Fraenkel, elección, Zorn y buen orden son formulaciones equivalentes. Esta equivalencia y su independencia respecto de los restantes axiomas, si éstos son consistentes, se toman aquí como hechos fundacionales. Se empleará Zorn en las demostraciones.

%%%%%%%%%%%%%%%%%%%%%%%%%%%%%%%%%%%%%%%%%%%%%%%%%%%%%%%%%%%%%%%%%%%%%%%%%%%%%%%
\subsection{Existencia de ideales maximales}
%%%%%%%%%%%%%%%%%%%%%%%%%%%%%%%%%%%%%%%%%%%%%%%%%%%%%%%%%%%%%%%%%%%%%%%%%%%%%%%
% P05.
\begin{prop}\label{prop:C17-maximales}
Todo ideal propio $I$ de un anillo conmutativo unitario $R$ está contenido en un ideal maximal.
\end{prop}
\begin{proof}
Ordenemos por inclusión el conjunto
\[
\mathcal I=\set{J\subsetneq R\given J\text{ es ideal e }I\subseteq J}.
\]
No es vacío, pues contiene $I$. Para una cadena no vacía $\mathcal C\subseteq\mathcal I$, sea $U=\bigcup_{J\in\mathcal C}J$. Contiene $I$. Si $a,b\in U$, pertenecen a dos miembros de la cadena y uno contiene al otro; así, ambos pertenecen a un mismo ideal y $a-b\in U$. También $ra\in U$ para todo $r\in R$. Luego $U$ es ideal.

Si $U=R$, entonces $1_R\in J$ para algún $J\in\mathcal C$, contradiciendo que $J$ es propio. Por tanto, $U\in\mathcal I$ y es cota superior de la cadena. La cadena vacía tiene como cota a $I$. Zorn da un elemento maximal $\mathfrak m$ de $\mathcal I$. Todo ideal propio que lo contenga también pertenece a $\mathcal I$, luego coincide con él. Así, $\mathfrak m$ es un ideal maximal de $R$.
\end{proof}
La unidad es esencial para concluir que la unión sigue siendo propia. En el anillo cero no hay ideales propios, de modo que el enunciado no afirma que tenga un maximal. Para un anillo no nulo, aplicarlo a $(0)$ sí asegura su existencia.

\begin{ejer}[La necesidad de la unidad]\label{ejer:C17-sin-unidad}
Dar un anillo sin unidad que no tenga ideales maximales.
\end{ejer}
\begin{sol}
Tomemos el grupo aditivo $\mathbb Q$ con multiplicación idénticamente cero. Es un anillo asociativo y conmutativo sin unidad: $e\cdot1=0\ne1$ para cualquier $e$. Todos sus subgrupos aditivos son ideales.

Si $M$ fuera un ideal maximal, el grupo abeliano no nulo $\mathbb Q/M$ no tendría subgrupos propios no triviales. Cualquier elemento no nulo lo generaría; sería cíclico, y su orden tendría que ser un primo $p$: un grupo cíclico infinito o de orden compuesto tiene subgrupos propios no triviales. Pero multiplicar por $p$ es sobreyectivo en $\mathbb Q/M$, porque $q+M=p(q/p+M)$. En un grupo de orden $p$ esa multiplicación es cero, una contradicción. No existe tal $M$.
\end{sol}

%%%%%%%%%%%%%%%%%%%%%%%%%%%%%%%%%%%%%%%%%%%%%%%%%%%%%%%%%%%%%%%%%%%%%%%%%%%%%%%
\section{Primitividad y contenido de polinomios}
%%%%%%%%%%%%%%%%%%%%%%%%%%%%%%%%%%%%%%%%%%%%%%%%%%%%%%%%%%%%%%%%%%%%%%%%%%%%%%%
%%%%%%%%%%%%%%%%%%%%%%%%%%%%%%%%%%%%%%%%%%%%%%%%%%%%%%%%%%%%%%%%%%%%%%%%%%%%%%%
\subsection{Extender un ideal al anillo de polinomios}
%%%%%%%%%%%%%%%%%%%%%%%%%%%%%%%%%%%%%%%%%%%%%%%%%%%%%%%%%%%%%%%%%%%%%%%%%%%%%%%
% P06: demostración manuscrita de la comprobación V.4.1.
Para un ideal $I$ de $R$, el ideal que genera en $R[x]$ es
\[
IR[x]=\set*{\sum_{i=0}^d a_ix^i\given d\geq0,\ a_i\in I}.
\]
En efecto, un polinomio con coeficientes en $I$ es suma de productos $a_ix^i$. A la inversa, multiplicar un elemento de $I$ por un polinomio da coeficientes en $I$, y las sumas conservan esa condición.

\begin{ejer}[Cocientar los coeficientes]\label{ejer:C17-coeficientes}
Probar el isomorfismo
\[
R[x]/IR[x]\cong(R/I)[x].
\]

\end{ejer}
\begin{sol}
Consideremos
\[
\varphi:R[x]\longrightarrow(R/I)[x],\qquad
\sum_i a_ix^i\longmapsto\sum_i(a_i+I)x^i.
\]
Preserva la suma por coeficientes y el producto porque la reducción de cada coeficiente de convolución satisface
\[
\left(\sum_{i+j=k}a_ib_j\right)+I
=\sum_{i+j=k}(a_i+I)(b_j+I).
\]
También preserva la unidad. Es sobreyectivo: se elige un representante de cada uno de los finitos coeficientes de un polinomio del destino. Su núcleo consiste exactamente en los polinomios cuyos coeficientes pertenecen a $I$, es decir, en $IR[x]$. El primer isomorfismo da el resultado.
\end{sol}
\begin{coro}\label{coro:C17-primo-extendido}
Si $I$ es primo en $R$, entonces $IR[x]$ es primo en $R[x]$.
\end{coro}
\begin{proof}
El cociente $(R/I)[x]$ es un dominio porque $R/I$ lo es. Aplicamos el isomorfismo anterior y el criterio de primalidad mediante cocientes.
\end{proof}

%%%%%%%%%%%%%%%%%%%%%%%%%%%%%%%%%%%%%%%%%%%%%%%%%%%%%%%%%%%%%%%%%%%%%%%%%%%%%%%
\subsection{Polinomios primitivos y muy primitivos}
%%%%%%%%%%%%%%%%%%%%%%%%%%%%%%%%%%%%%%%%%%%%%%%%%%%%%%%%%%%%%%%%%%%%%%%%%%%%%%%
% P07.
En el resto de la clase $R$ es conmutativo.
\begin{defi}
Un polinomio $f\in R[x]$ es \emph{muy primitivo} si $f\notin\mathfrak pR[x]$ para todo ideal primo $\mathfrak p$ de $R$. Es \emph{primitivo} si se exige lo mismo sólo para los ideales primos principales.
\end{defi}
\begin{lema}\label{lema:C17-producto-primitivo}
El producto $fg$ es primitivo si y sólo si ambos factores lo son.
\end{lema}
\begin{proof}
Para cada ideal primo principal $\mathfrak p$, el \cref{coro:C17-primo-extendido} muestra que $\mathfrak pR[x]$ es primo. Por tanto,
\[
fg\notin\mathfrak pR[x]
\quad\Longleftrightarrow\quad
f\notin\mathfrak pR[x]\ \text{y}\ g\notin\mathfrak pR[x].
\]
Exigirlo para todos esos ideales equivale a la primitividad de ambos factores.
\end{proof}
\begin{lema}\label{lema:C17-criterios-primitivo}
Para $f=a_0+\cdots+a_dx^d$ se cumplen las afirmaciones siguientes.
\begin{enumerate}[label=\textup{(\roman*)}]
\item $f$ es muy primitivo si y sólo si $(a_0,\ldots,a_d)=R$.
\item Si $R$ es DFU y $f\ne0$, es primitivo si y sólo si un máximo común divisor de sus coeficientes es una unidad.
\end{enumerate}
\end{lema}
\begin{proof}
Un ideal primo contiene $f$ coeficiente a coeficiente exactamente cuando contiene $(a_0,\ldots,a_d)$. Si este último es $R$, ningún primo puede contenerlo. Si es propio, la \cref{prop:C17-maximales} lo incluye en un maximal, que también es primo. Esto prueba (i).

Para (ii), si el máximo común divisor no es unidad, tiene un factor irreducible $q$. En un DFU, $(q)$ es primo y contiene todos los coeficientes, luego $f$ no es primitivo. Recíprocamente, si $f\in(q)R[x]$ para un ideal primo principal propio, su generador divide todos los coeficientes. Ese generador no puede ser cero, pues $f\ne0$, ni unidad, pues el ideal es propio. El máximo común divisor tampoco puede ser unidad.
\end{proof}

%%%%%%%%%%%%%%%%%%%%%%%%%%%%%%%%%%%%%%%%%%%%%%%%%%%%%%%%%%%%%%%%%%%%%%%%%%%%%%%
\subsection{Separar el contenido y multiplicarlo}
%%%%%%%%%%%%%%%%%%%%%%%%%%%%%%%%%%%%%%%%%%%%%%%%%%%%%%%%%%%%%%%%%%%%%%%%%%%%%%%
% P08--P10.
\begin{defi}
Si $R$ es DFU y $f\ne0$, su \emph{contenido} $c_f=\operatorname{cont}(f)$ es un máximo común divisor de sus coeficientes. Sólo está definido salvo unidades, pero el ideal principal $C_f=(c_f)$ está determinado de manera única.
\end{defi}
Así, $f$ es primitivo exactamente cuando $C_f=R$. Los productos de contenidos se expresarán mediante ideales: $(a)(b)=(ab)$.

\begin{ejem}
En $\mathbb Z[x]$, el polinomio $f=12x^2-18x+30$ tiene contenido $6$, de modo que
\[
f=6(2x^2-3x+5).
\]
No es primitivo, pero su factor $2x^2-3x+5$ sí lo es, pues $\mcd(2,-3,5)=1$.
\end{ejem}

\begin{lema}\label{lema:C17-parte-primitiva}
Si $R$ es DFU y $f\ne0$, existe $f_0\in R[x]$ primitivo con $f=c_ff_0$. Si $(f)=(c)(g)$ para $c\in R$ y $g$ primitivo, entonces $(c)=C_f$.
\end{lema}
\begin{proof}
Escribamos $a_i=c_fb_i$ para cada coeficiente y pongamos $f_0=\sum_i b_ix^i$. Si un irreducible $q$ dividiera todos los $b_i$, $c_fq$ dividiría todos los $a_i$; por la propiedad del máximo común divisor, $c_fq\mid c_f$. Cancelar $c_f\ne0$ haría a $q$ una unidad. Por tanto, los $b_i$ tienen máximo común divisor unidad y $f_0$ es primitivo.

Para la segunda parte, $(f)=(cg)$ implica $f=dcg$ y $cg=ef$ para ciertos polinomios $d,e$. Como $f\ne0$, cancelar en $f=def$ da $de=1$. En un anillo de polinomios sobre un dominio, las unidades son constantes: $\deg d+\deg e=0$. Así, $d\in R^\times$. Los coeficientes de $cg$ tienen máximo común divisor asociado a $c$, porque los de $g$ no tienen factor irreducible común. Multiplicar por $d$ no cambia su ideal de contenido. Luego $C_f=(c)$.
\end{proof}
\begin{obs}
El contenido se ha definido para polinomios no nulos. Las fórmulas siguientes se utilizan bajo esa hipótesis.
\end{obs}

\begin{prop}[Lema de Gauss]\label{prop:C17-Gauss}
Para $f,g\ne0$ en $R[x]$, con $R$ DFU, se tiene $C_{fg}=C_fC_g$.
\end{prop}
\begin{proof}
Escribimos $f=c_ff_0$ y $g=c_gg_0$, con $f_0,g_0$ primitivos. Su producto también es primitivo por el \cref{lema:C17-producto-primitivo}. La igualdad
\[
(fg)=(c_fc_g)(f_0g_0)
\]
y el \cref{lema:C17-parte-primitiva} dan $C_{fg}=(c_fc_g)=C_fC_g$.
\end{proof}
\begin{coro}\label{coro:C17-contenido-inclusion}
Si $f,g\ne0$ y $(f)\subseteq(g)$ en $R[x]$, entonces $C_f\subseteq C_g$.
\end{coro}
\begin{proof}
La inclusión dice $f=gh$ para algún $h\ne0$. Por Gauss, $C_f=C_gC_h\subseteq C_g$.
\end{proof}

\begin{obs}[Referencias de lectura]
Para Euclides y Zorn, véase \citet[cap.~V, 2.9--12 y 3.1/3/5]{aluffi2009}; para contenido y primitividad, V.4.1--9. Las equivalencias fundacionales se enuncian sin prueba. Los ejercicios corresponden a V.3.9 y V.4.1.
\end{obs}
